% Luchando los dias
%===========================================
% 16. Incorporación de ruido gaussiano moderado y Explicar Procedimientos para agregar ruido Gaussiano 
%===========================================

\begin{frame}{Incorporación de Ruido Gaussiano}

    \begin{columns}[T]
        %========================================
        % COLUMNA IZQUIERDA: Justificación Física
        %========================================
        \column{0.48\textwidth}
        \begin{block}{¿Por qué añadir ruido?}
            \scriptsize
            Para simular condiciones reales de medición:
            \begin{itemize}
                \item Errores instrumentales de los sensores.
                \item Variaciones e irregularidades del medio.
                \item Pequeñas perturbaciones externas.
            \end{itemize}
        \end{block}

        \vspace{0.2cm}

        \begin{block}{Naturaleza del Ruido ($\epsilon_i$)}
            \scriptsize
            \begin{itemize}
                \item \textbf{Distribución Normal:} Media cero ($\mu = 0$).
                \item \textbf{Imparcialidad:} No introduce tendencias; puede aumentar o disminuir el valor ideal.
            \end{itemize}
        \end{block}

        %========================================
        % COLUMNA DERECHA: Formulación Matemática
        %========================================
        \column{0.48\textwidth}
        \begin{exampleblock}{Modelado Matemático}
            \centering
            \footnotesize
            \textbf{Variación moderada (5\%)}
            \[ \text{Factor} = 0.05 \]
            
            \textbf{Desviación Estándar ($\sigma_i$)}
            \[ \sigma_i = 0.05 \cdot A_{ideal,i} \]
            
            \textbf{Generación de Ruido}
            \[ \epsilon_i \sim N(0, \sigma_i) \]
            
            \textbf{Amplitud Observada}
            \[ \mathbf{A_{obs,i} = A_{ideal,i} + \epsilon_i} \]
        \end{exampleblock}
    \end{columns}

    \vspace{0.4cm}
    \begin{center}
        \footnotesize 
        \textbf{El ruido transforma una simulación perfecta en un escenario de datos realistas.}
    \end{center}

\end{frame}